\documentclass[11pt,a4paper]{article}

\usepackage[utf8]{inputenc}
\usepackage[T1]{fontenc}
\usepackage[spanish]{babel}
\usepackage{amsmath,amssymb,mathtools}
\usepackage{geometry}
\usepackage{lmodern}
\usepackage{microtype}

\geometry{margin=2.4cm}
\setlength{\parindent}{0pt}
\setlength{\parskip}{0.45em}

\newcommand{\R}{\mathbb{R}}
\newcommand{\sen}{\operatorname{sen}}

\begin{document}

\begin{center}
{\LARGE\bfseries Análisis Matemático II}\\[0.4em]
{\Large Tema 4: Aplicaciones de la derivabilidad y diferenciabilidad en funciones de varias variables}\\[0.4em]
{\large Soluciones}
\end{center}

\section{Ejercicio 1}
\subsection{Solución}

Sea:

\[
f(x,y)=3x^2+y^2+2y.
\]

Los puntos críticos verifican:

\[
f_x=6x=0,
\qquad
f_y=2y+2=0.
\]

Por tanto:

\[
(x,y)=(0,-1).
\]

La matriz hessiana es:

\[
H_f=
\begin{pmatrix}
6&0\\
0&2
\end{pmatrix}.
\]

Como:

\[
|H_f|=12>0,
\qquad
f_{xx}=6>0,
\]

el punto es un mínimo relativo.

Además:

\[
f(0,-1)=-1.
\]

Por tanto:

\[
\boxed{
(0,-1)\text{ es un mínimo relativo, con }f(0,-1)=-1
}
\]

\section{Ejercicio 2}
\subsection{Solución}

Sea:

\[
f(x,y)=\frac8y+\frac yx+x,
\qquad
x>0,\ y>0.
\]

Las derivadas parciales son:

\[
f_x=1-\frac{y}{x^2},
\]

\[
f_y=-\frac8{y^2}+\frac1x.
\]

Igualando a cero:

\[
y=x^2,
\]

\[
y^2=8x.
\]

Sustituyendo:

\[
x^4=8x
\]

y, como \(x>0\):

\[
x=2,
\qquad
y=4.
\]

La matriz hessiana es:

\[
H_f=
\begin{pmatrix}
\dfrac{2y}{x^3}&-\dfrac1{x^2}\\[2mm]
-\dfrac1{x^2}&\dfrac{16}{y^3}
\end{pmatrix}.
\]

En \((2,4)\):

\[
H_f(2,4)=
\begin{pmatrix}
1&-\frac14\\[1mm]
-\frac14&\frac14
\end{pmatrix}.
\]

Entonces:

\[
|H_f(2,4)|
=
\frac3{16}>0,
\qquad
f_{xx}(2,4)=1>0.
\]

Por tanto:

\[
\boxed{
(2,4)\text{ es un mínimo relativo}
}
\]

y:

\[
\boxed{
f(2,4)=6
}
\]

\section{Ejercicio 3}
\subsection{Solución}

Sea:

\[
f(x,y)
=
x^2(x^2-2)+y^2(y^2-2)+4xy.
\]

Las derivadas parciales son:

\[
f_x
=
4(x^3-x+y),
\]

\[
f_y
=
4(y^3-y+x).
\]

Los puntos críticos satisfacen:

\[
x^3-x+y=0,
\]

\[
y^3-y+x=0.
\]

Restando:

\[
(x-y)(x^2+xy+y^2-2)=0.
\]

Sumando:

\[
(x+y)(x^2-xy+y^2)=0.
\]

Se obtienen:

\[
\boxed{
(0,0),
\quad
(\sqrt2,-\sqrt2),
\quad
(-\sqrt2,\sqrt2)
}
\]

La matriz hessiana es:

\[
H_f=
\begin{pmatrix}
12x^2-4&4\\
4&12y^2-4
\end{pmatrix}.
\]

En:

\[
(\pm\sqrt2,\mp\sqrt2),
\]

\[
H_f=
\begin{pmatrix}
20&4\\
4&20
\end{pmatrix}.
\]

Como:

\[
|H_f|=384>0,
\qquad
f_{xx}=20>0,
\]

ambos son mínimos relativos.

Además:

\[
f(\sqrt2,-\sqrt2)
=
f(-\sqrt2,\sqrt2)
=
-8.
\]

En \((0,0)\):

\[
|H_f(0,0)|=0.
\]

Estudiamos dos direcciones:

\[
f(x,x)=2x^4\geq0,
\]

\[
f(x,-x)=2x^4-8x^2<0
\]

para \(x\) suficientemente próximo a \(0\), \(x\neq0\).

Por tanto:

\[
\boxed{
(0,0)\text{ es un punto de silla}
}
\]

\[
\boxed{
(\sqrt2,-\sqrt2)
\text{ y }
(-\sqrt2,\sqrt2)
\text{ son mínimos relativos de valor }-8
}
\]

\section{Ejercicio 4}
\subsection{Solución}

Sea:

\[
f(x,y)
=
x^4+y^4+\frac1{x^4y^4},
\qquad
x\neq0,\ y\neq0.
\]

Las derivadas parciales son:

\[
f_x
=
4x^3-\frac4{x^5y^4},
\]

\[
f_y
=
4y^3-\frac4{x^4y^5}.
\]

Igualando a cero:

\[
x^8y^4=1,
\]

\[
x^4y^8=1.
\]

Dividiendo ambas ecuaciones:

\[
x^4=y^4.
\]

Por tanto:

\[
|x|=|y|.
\]

Sustituyendo:

\[
x^{12}=1.
\]

Así, los puntos críticos son:

\[
\boxed{
(1,1),\ (1,-1),\ (-1,1),\ (-1,-1)
}
\]

Las derivadas de segundo orden son:

\[
f_{xx}
=
12x^2+\frac{20}{x^6y^4},
\]

\[
f_{yy}
=
12y^2+\frac{20}{x^4y^6},
\]

\[
f_{xy}
=
\frac{16}{x^5y^5}.
\]

En cualquiera de los puntos críticos:

\[
f_{xx}=f_{yy}=32,
\]

\[
|H_f|
=
32^2-16^2
=
768>0.
\]

Por tanto, todos son mínimos relativos.

Además:

\[
f(\pm1,\pm1)=3.
\]

Así:

\[
\boxed{
(\pm1,\pm1)\text{ son mínimos relativos de valor }3
}
\]

\section{Ejercicio 5}
\subsection{Solución}

Sea:

\[
f(x,y)
=
xe^y+ye^x.
\]

Las derivadas parciales son:

\[
f_x=e^y+ye^x,
\]

\[
f_y=xe^y+e^x.
\]

En un punto crítico:

\[
e^y=-ye^x,
\]

\[
e^x=-xe^y.
\]

Multiplicando:

\[
xy=1.
\]

Además \(x<0\) e \(y<0\), por lo que:

\[
y=\frac1x.
\]

La única solución es:

\[
\boxed{
(x,y)=(-1,-1)
}
\]

La matriz hessiana es:

\[
H_f=
\begin{pmatrix}
ye^x&e^x+e^y\\
e^x+e^y&xe^y
\end{pmatrix}.
\]

En \((-1,-1)\):

\[
H_f(-1,-1)
=
\frac1e
\begin{pmatrix}
-1&2\\
2&-1
\end{pmatrix}.
\]

Entonces:

\[
|H_f(-1,-1)|
=
-\frac3{e^2}<0.
\]

Por tanto:

\[
\boxed{
(-1,-1)\text{ es un punto de silla}
}
\]

y la función no tiene máximos ni mínimos relativos.

\section{Ejercicio 6}
\subsection{Solución}

Queremos extremar:

\[
f(x,y)=3x+y
\]

sujeto a:

\[
x^2+y^2=1000.
\]

Definimos:

\[
L(x,y,\lambda)
=
3x+y-\lambda(x^2+y^2-1000).
\]

Las ecuaciones de Lagrange son:

\[
3=2\lambda x,
\]

\[
1=2\lambda y,
\]

\[
x^2+y^2=1000.
\]

Dividiendo:

\[
x=3y.
\]

Entonces:

\[
10y^2=1000
\]

y:

\[
y=\pm10.
\]

Por tanto:

\[
(x,y)=(30,10),
\qquad
(x,y)=(-30,-10).
\]

Los valores son:

\[
f(30,10)=100,
\]

\[
f(-30,-10)=-100.
\]

Así:

\[
\boxed{
f_{\max}=100
\text{ en }(30,10)
}
\]

\[
\boxed{
f_{\min}=-100
\text{ en }(-30,-10)
}
\]

\section{Ejercicio 7}
\subsection{Solución}

Queremos maximizar:

\[
f(x,y)=xy
\]

sujeto a:

\[
2x+2y=a.
\]

La restricción equivale a:

\[
x+y=\frac a2.
\]

Definimos:

\[
L(x,y,\lambda)
=
xy-\lambda\left(2x+2y-a\right).
\]

Entonces:

\[
y-2\lambda=0,
\]

\[
x-2\lambda=0.
\]

Por tanto:

\[
x=y.
\]

Usando la restricción:

\[
4x=a.
\]

Así:

\[
x=y=\frac a4.
\]

Por tanto:

\[
\boxed{
f_{\max}
=
\frac{a^2}{16}
}
\]

alcanzado en:

\[
\boxed{
\left(\frac a4,\frac a4\right)
}
\]

\section{Ejercicio 8}
\subsection{Solución}

La distancia desde:

\[
P=(1,-2)
\]

hasta un punto \((x,y)\) de la recta:

\[
x+5y-10=0
\]

es:

\[
d^2
=
(x-1)^2+(y+2)^2.
\]

Aplicamos Lagrange a:

\[
f(x,y)
=
(x-1)^2+(y+2)^2
\]

con la restricción:

\[
g(x,y)=x+5y-10=0.
\]

Se obtiene:

\[
2(x-1)=\lambda,
\]

\[
2(y+2)=5\lambda.
\]

Resolviendo junto con la restricción:

\[
\boxed{
(x,y)
=
\left(
\frac{45}{26},
\frac{43}{26}
\right)
}
\]

La distancia mínima es:

\[
d_{\min}
=
\sqrt{
\left(\frac{19}{26}\right)^2
+
\left(\frac{95}{26}\right)^2
}.
\]

Por tanto:

\[
\boxed{
d_{\min}
=
\frac{19}{\sqrt{26}}
}
\]

\section{Ejercicio 9}
\subsection{Solución}

Para un punto \((x,y)\) de la elipse:

\[
x^2+4y^2=4,
\]

la distancia a la recta:

\[
x+y=4
\]

es:

\[
d
=
\frac{|x+y-4|}{\sqrt2}.
\]

Como en la elipse:

\[
x+y<4,
\]

minimizar la distancia equivale a maximizar:

\[
x+y.
\]

Aplicamos Lagrange:

\[
1=2\lambda x,
\]

\[
1=8\lambda y.
\]

Por tanto:

\[
x=4y.
\]

Usando la restricción:

\[
16y^2+4y^2=4.
\]

Para el máximo de \(x+y\):

\[
y=\frac1{\sqrt5},
\qquad
x=\frac4{\sqrt5}.
\]

Así:

\[
x+y=\sqrt5.
\]

Por tanto:

\[
\boxed{
d_{\min}
=
\frac{4-\sqrt5}{\sqrt2}
}
\]

\section{Ejercicio 10}
\subsection{Solución}

Sea:

\[
f(x,y,z)=x-y+z
\]

sujeto a:

\[
x^2+y^2+z^2=2.
\]

Definimos:

\[
L
=
x-y+z
-
\lambda(x^2+y^2+z^2-2).
\]

Las ecuaciones son:

\[
1=2\lambda x,
\]

\[
-1=2\lambda y,
\]

\[
1=2\lambda z.
\]

Por tanto:

\[
x=z=-y.
\]

Usando la restricción:

\[
3x^2=2.
\]

Así:

\[
x=\pm\sqrt{\frac23}.
\]

Por tanto:

\[
\boxed{
f_{\max}=\sqrt6
}
\]

en:

\[
\boxed{
\sqrt{\frac23}(1,-1,1)
}
\]

y:

\[
\boxed{
f_{\min}=-\sqrt6
}
\]

en:

\[
\boxed{
-\sqrt{\frac23}(1,-1,1)
}
\]

\section{Ejercicio 11}
\subsection{Solución}

Sea:

\[
P=(2,1,-2).
\]

Su distancia al origen es:

\[
\|P\|
=
\sqrt{4+1+4}
=
3.
\]

La esfera:

\[
x^2+y^2+z^2=1
\]

tiene radio \(1\).

La distancia mínima desde \(P\) a la esfera es:

\[
3-1=2,
\]

y la máxima:

\[
3+1=4.
\]

Por tanto:

\[
\boxed{
d_{\min}=2
}
\]

\[
\boxed{
d_{\max}=4
}
\]

Los puntos correspondientes son:

\[
\boxed{
\left(
\frac23,\frac13,-\frac23
\right)
}
\]

y:

\[
\boxed{
\left(
-\frac23,-\frac13,\frac23
\right)
}
\]

respectivamente.

\section{Ejercicio 12}
\subsection{Solución}

Sean \(x,y,z>0\) las aristas del paralelepípedo.

Su volumen es:

\[
V=xyz.
\]

La superficie total es:

\[
2(xy+xz+yz)=S.
\]

Aplicamos Lagrange:

\[
L
=
xyz
-
\lambda\left(
2xy+2xz+2yz-S
\right).
\]

Las ecuaciones son:

\[
yz
=
2\lambda(y+z),
\]

\[
xz
=
2\lambda(x+z),
\]

\[
xy
=
2\lambda(x+y).
\]

Comparando se obtiene:

\[
x=y=z.
\]

Sea:

\[
x=y=z=a.
\]

Entonces:

\[
6a^2=S.
\]

Por tanto:

\[
a=\sqrt{\frac S6}.
\]

Así, el paralelepípedo de mayor volumen es un cubo y:

\[
\boxed{
x=y=z=\sqrt{\frac S6}
}
\]

\[
\boxed{
V_{\max}
=
\left(
\frac S6
\right)^{3/2}
}
\]

\section{Ejercicio 13}
\subsection{Solución}

Sea:

\[
f(x,y)=3xy
\]

en:

\[
x^2+y^2\leq3.
\]

En el interior:

\[
f_x=3y,
\qquad
f_y=3x.
\]

El único punto crítico es:

\[
(0,0),
\]

donde:

\[
f(0,0)=0.
\]

En la frontera:

\[
x^2+y^2=3.
\]

Aplicamos Lagrange:

\[
3y=2\lambda x,
\]

\[
3x=2\lambda y.
\]

Para los extremos no nulos:

\[
x^2=y^2.
\]

El máximo se obtiene cuando:

\[
x=y.
\]

Entonces:

\[
2x^2=3.
\]

Por tanto:

\[
x=y=\pm\sqrt{\frac32}.
\]

Así:

\[
\boxed{
f_{\max}
=
\frac92
}
\]

alcanzado en:

\[
\boxed{
\left(
\sqrt{\frac32},
\sqrt{\frac32}
\right),
\qquad
\left(
-\sqrt{\frac32},
-\sqrt{\frac32}
\right)
}
\]

\section{Ejercicio 14}
\subsection{Solución}

Sea:

\[
f(x,y,z)
=
-x+y-3z
\]

en:

\[
x^2+2y^2+z^2\leq7.
\]

Como \(f\) es lineal, los extremos absolutos se alcanzan en la frontera:

\[
x^2+2y^2+z^2=7.
\]

Aplicamos Lagrange:

\[
-1=2\lambda x,
\]

\[
1=4\lambda y,
\]

\[
-3=2\lambda z.
\]

Así:

\[
x=-\frac1{2\lambda},
\qquad
y=\frac1{4\lambda},
\qquad
z=-\frac3{2\lambda}.
\]

Sustituyendo en la restricción:

\[
\frac{21}{8\lambda^2}=7.
\]

Por tanto:

\[
\lambda^2=\frac38.
\]

Los puntos extremos son:

\[
P_{\max}
=
\sqrt{\frac23}
\left(
-1,\frac12,-3
\right),
\]

\[
P_{\min}
=
-\sqrt{\frac23}
\left(
-1,\frac12,-3
\right).
\]

Los valores son:

\[
\boxed{
f_{\max}
=
\frac{7\sqrt6}{2}
}
\]

\[
\boxed{
f_{\min}
=
-\frac{7\sqrt6}{2}
}
\]

Estos extremos son absolutos y también relativos respecto del conjunto \(S\).

\section{Ejercicio 15}
\subsection{Solución}

Sea:

\[
f(x,y)
=
8x^3-24xy+y^3
\]

en el triángulo:

\[
x\geq0,
\qquad
y\geq0,
\qquad
x+y\leq1.
\]

En el interior:

\[
f_x=24x^2-24y,
\]

\[
f_y=-24x+3y^2.
\]

Las ecuaciones:

\[
y=x^2,
\qquad
y^2=8x
\]

no producen ningún punto crítico interior.

Estudiamos la frontera.

Para:

\[
x=0,
\]

\[
f(0,y)=y^3,
\qquad
0\leq y\leq1.
\]

Para:

\[
y=0,
\]

\[
f(x,0)=8x^3,
\qquad
0\leq x\leq1.
\]

Para:

\[
x+y=1,
\qquad
y=1-x,
\]

se tiene:

\[
f(x,1-x)
=
7x^3+27x^2-27x+1.
\]

Derivando:

\[
\frac{d}{dx}f(x,1-x)
=
3(x+3)(7x-3).
\]

El punto relevante es:

\[
x=\frac37,
\qquad
y=\frac47.
\]

Su valor es:

\[
f\left(
\frac37,\frac47
\right)
=
-\frac{248}{49}.
\]

En los vértices:

\[
f(0,0)=0,
\]

\[
f(0,1)=1,
\]

\[
f(1,0)=8.
\]

Por tanto:

\[
\boxed{
f_{\max}=8
\text{ en }(1,0)
}
\]

\[
\boxed{
f_{\min}
=
-\frac{248}{49}
\text{ en }
\left(
\frac37,\frac47
\right)
}
\]

\section{Ejercicio 16}
\subsection{Solución}

Sea:

\[
f(x,y)
=
e^{x^2-2x+y^2}.
\]

Escribimos:

\[
x^2-2x+y^2
=
(x-1)^2+y^2-1.
\]

El recinto es:

\[
(x-1)^2+2y^2\leq1.
\]

Como la función exponencial es creciente, basta extremar:

\[
q(x,y)
=
(x-1)^2+y^2.
\]

En el recinto:

\[
q\geq0,
\]

y el mínimo se alcanza en:

\[
(x,y)=(1,0).
\]

Por tanto:

\[
f_{\min}
=
e^{-1}.
\]

Para el máximo:

\[
(x-1)^2+y^2
\leq
(x-1)^2+2y^2
\leq1.
\]

La igualdad se alcanza cuando:

\[
y=0,
\qquad
(x-1)^2=1.
\]

Es decir:

\[
(x,y)=(0,0),
\qquad
(x,y)=(2,0).
\]

Así:

\[
\boxed{
f_{\min}
=
e^{-1}
\text{ en }(1,0)
}
\]

\[
\boxed{
f_{\max}
=
1
\text{ en }(0,0)\text{ y }(2,0)
}
\]

\section{Ejercicio 17}
\subsection{Solución}

Queremos extremar:

\[
f(x,y,z)=z
\]

sujeto a:

\[
x+y+z=12,
\]

\[
z=x^2+y^2.
\]

De la primera condición:

\[
x+y=12-z.
\]

Además:

\[
x^2+y^2=z.
\]

Por Cauchy--Schwarz:

\[
(x+y)^2
\leq
2(x^2+y^2).
\]

Por tanto:

\[
(12-z)^2\leq2z.
\]

Así:

\[
z^2-26z+144\leq0.
\]

Factorizando:

\[
(z-8)(z-18)\leq0.
\]

Luego:

\[
8\leq z\leq18.
\]

Los extremos se alcanzan cuando:

\[
x=y.
\]

Para \(z=8\):

\[
x=y=2.
\]

Para \(z=18\):

\[
x=y=-3.
\]

Por tanto:

\[
\boxed{
z_{\min}=8
\text{ en }(2,2,8)
}
\]

\[
\boxed{
z_{\max}=18
\text{ en }(-3,-3,18)
}
\]

\section{Ejercicio 18}
\subsection{Solución}

Las restricciones son:

\[
x+y+z=40,
\]

\[
z=x+y.
\]

Sustituyendo:

\[
2z=40.
\]

Por tanto:

\[
z=20,
\qquad
x+y=20.
\]

La función queda:

\[
f(x,y,z)
=
20xy.
\]

Para una suma fija:

\[
x+y=20,
\]

el producto \(xy\) es máximo cuando:

\[
x=y=10.
\]

Por tanto:

\[
\boxed{
f_{\max}
=
10\cdot10\cdot20
=
2000
}
\]

alcanzado en:

\[
\boxed{
(10,10,20)
}
\]

\section{Ejercicio 19}
\subsection{Solución}

Sea:

\[
f(x,y)=e^{2x+y}.
\]

En \((0,0)\):

\[
f(0,0)=1.
\]

Las derivadas de primer orden son:

\[
f_x=2e^{2x+y},
\qquad
f_y=e^{2x+y}.
\]

Por tanto:

\[
f_x(0,0)=2,
\qquad
f_y(0,0)=1.
\]

El polinomio de Maclaurin de primer grado es:

\[
\boxed{
P_1(x,y)
=
1+2x+y
}
\]

Las derivadas de segundo orden son:

\[
f_{xx}=4e^{2x+y},
\]

\[
f_{xy}=f_{yx}=2e^{2x+y},
\]

\[
f_{yy}=e^{2x+y}.
\]

En el origen:

\[
f_{xx}=4,
\qquad
f_{xy}=f_{yx}=2,
\qquad
f_{yy}=1.
\]

Por tanto:

\[
\boxed{
P_2(x,y)
=
1+2x+y
+
2x^2+2xy+\frac12y^2
}
\]

\section{Ejercicio 20}
\subsection{Solución}

Sea:

\[
f(x,y)
=
x\cos y+y\sen x.
\]

En el origen:

\[
f(0,0)=0.
\]

Las derivadas de primer orden son:

\[
f_x
=
\cos y+y\cos x,
\]

\[
f_y
=
-x\sen y+\sen x.
\]

Así:

\[
f_x(0,0)=1,
\qquad
f_y(0,0)=0.
\]

Las derivadas de segundo orden son:

\[
f_{xx}
=
-y\sen x,
\]

\[
f_{xy}
=
-\sen y+\cos x,
\]

\[
f_{yx}
=
-\sen y+\cos x,
\]

\[
f_{yy}
=
-x\cos y.
\]

En \((0,0)\):

\[
f_{xx}=0,
\qquad
f_{xy}=f_{yx}=1,
\qquad
f_{yy}=0.
\]

Por tanto:

\[
\boxed{
P_2(x,y)
=
x+xy
}
\]

\section{Ejercicio 21}
\subsection{Solución}

Sea:

\[
f(x,y)
=
\ln\left(
\frac{1+x}{1+y}
\right).
\]

Escribimos:

\[
f(x,y)
=
\ln(1+x)-\ln(1+y).
\]

Utilizando:

\[
\ln(1+t)
=
t-\frac{t^2}{2}+O(t^3),
\]

se obtiene:

\[
f(x,y)
=
x-\frac{x^2}{2}
-
y+\frac{y^2}{2}
+
O\left(\|(x,y)\|^3\right).
\]

Por tanto:

\[
\boxed{
P_2(x,y)
=
x-y-\frac{x^2}{2}+\frac{y^2}{2}
}
\]

\section{Ejercicio 22}
\subsection{Solución}

Sea:

\[
f(x,y)
=
x^3+y^3-x^2y+xy-1.
\]

Queremos expresarla en potencias de:

\[
x+1
\qquad\text{e}\qquad
y-2.
\]

Tomamos:

\[
h=x+1,
\qquad
k=y-2.
\]

Entonces:

\[
x=h-1,
\qquad
y=k+2.
\]

Sustituyendo y desarrollando:

\[
f
=
h^3-h^2k-5h^2+3hk+9h
+k^3+6k^2+10k+2.
\]

Por tanto:

\[
\boxed{
\begin{aligned}
f(x,y)
={}&
(x+1)^3
-(x+1)^2(y-2)
-5(x+1)^2\\
&+
3(x+1)(y-2)
+
9(x+1)
+
(y-2)^3\\
&+
6(y-2)^2
+
10(y-2)
+
2.
\end{aligned}
}
\]

\section{Ejercicio 23}
\subsection{Solución}

Sea:

\[
f(x,y)=x^y.
\]

Desarrollamos alrededor de:

\[
(1,0).
\]

En dicho punto:

\[
f(1,0)=1.
\]

Además:

\[
f_x
=
yx^{y-1},
\]

\[
f_y
=
x^y\ln x.
\]

Por tanto:

\[
f_x(1,0)=0,
\qquad
f_y(1,0)=0.
\]

Las derivadas de segundo orden son:

\[
f_{xx}
=
y(y-1)x^{y-2},
\]

\[
f_{xy}
=
x^{y-1}
+
yx^{y-1}\ln x,
\]

\[
f_{yy}
=
x^y(\ln x)^2.
\]

En \((1,0)\):

\[
f_{xx}=0,
\qquad
f_{xy}=f_{yx}=1,
\qquad
f_{yy}=0.
\]

Así:

\[
\boxed{
P_2(x,y)
=
1+(x-1)y
}
\]

Para:

\[
x=1.2,
\qquad
y=0.1,
\]

se obtiene:

\[
(1.2)^{0.1}
\approx
1+(0.2)(0.1).
\]

Por tanto:

\[
\boxed{
(1.2)^{0.1}
\approx
1.02
}
\]

\section{Ejercicio 24}
\subsection{Solución}

Tomamos:

\[
f(x,y)
=
\sqrt{x^2+y^3}
\]

y desarrollamos alrededor de:

\[
(1,2).
\]

Se tiene:

\[
f(1,2)=3.
\]

Las derivadas en \((1,2)\) son:

\[
f_x(1,2)=\frac13,
\]

\[
f_y(1,2)=2,
\]

\[
f_{xx}(1,2)=\frac8{27},
\]

\[
f_{xy}(1,2)=f_{yx}(1,2)=-\frac29,
\]

\[
f_{yy}(1,2)=\frac23.
\]

Por tanto:

\[
\boxed{
\begin{aligned}
P_2(x,y)
={}&
3
+
\frac13(x-1)
+
2(y-2)\\
&+
\frac4{27}(x-1)^2
-
\frac29(x-1)(y-2)
+
\frac13(y-2)^2.
\end{aligned}
}
\]

Para:

\[
x=1.02,
\qquad
y=1.97,
\]

\[
x-1=0.02,
\qquad
y-2=-0.03.
\]

Así:

\[
P_2(1.02,1.97)
\approx
2.947159.
\]

Por tanto:

\[
\boxed{
\sqrt{(1.02)^2+(1.97)^3}
\approx
2.94716
}
\]

\section{Ejercicio 25}
\subsection{Solución}

Sea:

\[
f(x,y,z)
=
e^{a(x+y+z)}.
\]

En el origen:

\[
f(0,0,0)=1.
\]

Como:

\[
e^t
=
1+t+\frac{t^2}{2}+O(t^3),
\]

tomando:

\[
t=a(x+y+z),
\]

se obtiene:

\[
\boxed{
P_2(x,y,z)
=
1
+
a(x+y+z)
+
\frac{a^2}{2}(x+y+z)^2
}
\]

\section{Ejercicio 26}
\subsection{Solución}

Sea:

\[
f(x,y)
=
\ln\left(\frac xy\right)
=
\ln x-\ln y.
\]

Desarrollamos alrededor de:

\[
(1,1).
\]

Tomamos:

\[
h=x-1,
\qquad
k=y-1.
\]

Entonces:

\[
f(x,y)
=
\ln(1+h)-\ln(1+k).
\]

Usando:

\[
\ln(1+t)
=
t-\frac{t^2}{2}+\frac{t^3}{3}+O(t^4),
\]

se obtiene:

\[
\boxed{
\begin{aligned}
P_3(x,y)
={}&
(x-1)-(y-1)\\
&-
\frac12
\left[
(x-1)^2-(y-1)^2
\right]\\
&+
\frac13
\left[
(x-1)^3-(y-1)^3
\right].
\end{aligned}
}
\]

\end{document}
